\subsection{Alignment results}
\IfFileExists{paper_ux/tables/alignment_main.tex}{\begin{table}[t]
\centering
\small
\caption{Proxy alignment summary (top-$k$ Jaccard and weighted-Jaccard). For embedding alignment, we embed a subsample of $n=200$ texts for local compute; on large app-review corpora, we additionally cap the proxy slice at $n=10{,}000$ texts.}
\label{tab:alignment-main}
\begin{tabular}{llrrr}
\toprule
Proxy & Method & $k$ & $J_k$ & $W$ \\
\midrule
\DatasetTinder & \MethodLex & 8 & 0.6 & 0.001 \\
\DatasetTinder & \MethodEmbed & 8 & 0.6 & 0.087 \\
\DatasetGojek & \MethodLex & 8 & 0.6 & 0 \\
\DatasetGojek & \MethodEmbed & 8 & 1 & 0.128 \\
\DatasetTwitter & \MethodLex & 8 & 0.6 & 0 \\
\DatasetTwitter & \MethodEmbed & 8 & 0.778 & 0.123 \\
\DatasetAmazon & \MethodLex & 8 & 0.6 & 0.009 \\
\DatasetAmazon & \MethodTFIDF & 8 & 0.6 & 0.001 \\
\DatasetAmazon & \MethodEmbed & 8 & 0.778 & 0.119 \\
\DatasetAmazonAnn & \MethodLex & 8 & 0.6 & 0.013 \\
\DatasetAmazonAnn & \MethodTFIDF & 8 & 0.778 & 0.003 \\
\DatasetAmazonAnn & \MethodEmbed & 8 & 0.778 & 0.061 \\
\bottomrule
\end{tabular}
\end{table}
}{}
\UXResolve{\uxalignpdf}{figures/alignment_plot.pdf}
\ifx\uxalignpdf\empty
  \UXResolve{\uxaligntex}{figures/alignment_plot.tex}
  \ifx\uxaligntex\empty\else
  \begin{figure}[t]
    \centering
    \resizebox{0.98\linewidth}{!}{\input{\uxaligntex}}
    \caption{Proxy alignment by dataset and method (weighted-Jaccard; higher is better).}
    \label{fig:alignment}
  \end{figure}
  \fi
\else
  \begin{figure}[t]
    \centering
    \includegraphics[width=0.98\linewidth]{\uxalignpdf}
    \caption{Proxy alignment by dataset and method (weighted-Jaccard; higher is better).}
    \label{fig:alignment}
  \end{figure}
\fi

Across our main proxy corpora (app reviews and support tweets), \MethodEmbed{} produces higher weighted-Jaccard overlap than lexical or TF-IDF baselines (Figure~\ref{fig:alignment}), suggesting that
multilingual semantic similarity is better aligned with our friction taxonomy than keyword matching.
However, absolute overlaps remain modest (Table~\ref{tab:alignment-main}), which is expected: proxy corpora are noisy mixtures of domains and user populations,
while simulation outputs are constrained by our prompt and taxonomy.
In OSS issue proxies, we observe exceptions where lexical overlap can exceed embedding overlap (see discussion below), so we treat the embedding advantage as a tendency rather than a universal guarantee.

\paragraph{Interpreting overlap (example).}
\UXResolve{\uxdistpdf}{figures/distribution_amazon_ann.pdf}
\ifx\uxdistpdf\empty
  \UXResolve{\uxdisttex}{figures/distribution_amazon_ann.tex}
  \ifx\uxdisttex\empty\else
  \begin{figure}[t]
    \centering
    \input{\uxdisttex}
    \caption{Friction-theme distributions for simulations vs proxy corpus (Amazon annotated; lexical mapping; top 8 categories). The simulation bars are near-zero because lexical keyword matching fails to match the structured taxonomy-prefixed outputs produced by the v2 prompt---this is an intentional diagnostic: it shows that lexical alignment is insufficient for simulation outputs, motivating the embedding baseline.}
    \label{fig:distribution-example}
  \end{figure}
  \fi
\else
  \begin{figure}[t]
    \centering
    \includegraphics[width=0.98\linewidth]{\uxdistpdf}
    \caption{Friction-theme distributions for simulations vs proxy corpus (Amazon annotated; lexical mapping; top 8 categories). The simulation bars are near-zero because lexical keyword matching fails to match the structured taxonomy-prefixed outputs produced by the v2 prompt---this is an intentional diagnostic: it shows that lexical alignment is insufficient for simulation outputs, motivating the embedding baseline.}
    \label{fig:distribution-example}
  \end{figure}
\fi
The goal is not to achieve perfect distributional match; rather, Figure~\ref{fig:distribution-example} helps diagnose \emph{where} simulations and proxies diverge.
Note that the aggregate $W$ values in Table~\ref{tab:alignment-main} are computed from category-level counts accumulated across the full proxy corpus under a shared taxonomy, and may differ from the per-run distribution bars shown here.
The sensitivity rows in Table~\ref{tab:alignment-sensitivity} were executed on a separate simulation pool dedicated to sweeping $k$ and threshold settings, which is why their absolute $W$ values for the same proxy/method pair (e.g., \DatasetAmazonAnn{} lexical, $W=0.017$) can differ slightly from the canonical entry in Table~\ref{tab:alignment-main} ($W=0.013$); we treat Table~\ref{tab:alignment-main} as the canonical reference and use the sensitivity rows only for relative comparisons across $k$ and thresholds.

\paragraph{Key takeaways.}
\begin{itemize}
  \item \MethodEmbed{} typically yields higher $W$ than \MethodLex{} on our primary proxies, making it a strong default under a fixed taxonomy.
  \item $J_k$ is useful for quick inspection but can mislead under thresholding or large $k$; we treat $W$ as the primary signal.
  \item The goal is not to ``prove UX quality'' but to enable early-stage iteration with auditable artifacts and bounded claims.
\end{itemize}

\paragraph{What a consumer can do with these results.}
In practice, the main value is comparative and diagnostic: a team can compare prompt/taxonomy versions (or model settings)
and see whether friction distributions move in the expected direction across proxies, rather than relying on anecdotes.
The same artifacts also support a pragmatic workflow: pick the top-$k$ friction themes to tackle, ship a change, and rerun the pipeline to check
whether the proxy-aligned distribution shifts (and whether the shift is stable under bootstrap intervals).

\paragraph{B2B proxy (OSS issues).}
On GitHub issues from observability repositories, embedding-based alignment tends to outperform lexical matching (Table~\ref{tab:alignment-oss}).
\IfFileExists{paper_ux/tables/alignment_oss.tex}{\begin{table}[t]
\centering
\small
\caption{OSS issues proxy alignment (selected repositories).}
\label{tab:alignment-oss}
\begin{tabular}{lrr}
\toprule
Repository & $W$ (\MethodLex) & $W$ (\MethodEmbed) \\
\midrule
Grafana & 0.12 & 0.185 \\
Jaeger & 0.078 & 0.1 \\
OpenTelemetry Collector & 0.1 & 0.157 \\
OpenTelemetry JS & 0.101 & 0.103 \\
Prometheus & 0.199 & 0.254 \\
Sentry & 0.138 & 0.072 \\
Zipkin & 0.073 & 0.208 \\
\bottomrule
\end{tabular}
\end{table}
}{}
While these proxies are not controlled UX studies, they broaden the evidence beyond app reviews and help stress-test whether our taxonomy and prompts
generalize to B2B tooling discussions.
We also observe failure modes: for some repositories (e.g., Sentry in our selection), lexical overlap can exceed embedding overlap,
suggesting that issue text may be dominated by product-specific vocabulary that semantic embeddings smooth over.

We also observe that top-$k$ Jaccard can be misleading when $k$ is too large (Table~\ref{tab:alignment-sensitivity}).
In particular, lexical $J_k$ reaches 1.0 at $k=10$ on the Amazon annotated proxy while weighted-Jaccard remains unchanged,
supporting our choice to emphasize $W$ for distributional comparisons and to keep $k$ smaller than the taxonomy size.

\subsection{Bootstrap confidence intervals}
\label{sec:results-bootstrap}
\IfFileExists{paper_ux/tables/bootstrap_ci.tex}{\begin{table}[t]
\centering
\small
\caption{Bootstrap confidence intervals for alignment metrics (Amazon low-rated; lexical).}
\label{tab:bootstrap-ci}
\begin{tabular}{lrr}
\toprule
Metric & Mean & 95\% CI \\
\midrule
$J_k$ & 0.601 & [0.6, 0.6] \\
$W$ & 0.011 & [0.01, 0.011] \\
\bottomrule
\end{tabular}
\end{table}
}{}
\UXResolve{\uxbootpdf}{figures/bootstrap_ci_plot.pdf}
\ifx\uxbootpdf\empty
  \UXResolve{\uxboottex}{figures/bootstrap_ci_plot.tex}
  \ifx\uxboottex\empty\else
  \begin{figure}[t]
    \centering
    \input{\uxboottex}
    \caption{Bootstrap confidence interval for weighted-Jaccard on Amazon low-rated (lexical).}
    \label{fig:bootstrap-ci}
  \end{figure}
  \fi
\else
  \begin{figure}[t]
    \centering
    \includegraphics[width=0.72\linewidth]{\uxbootpdf}
    \caption{Bootstrap confidence interval for weighted-Jaccard on Amazon low-rated (lexical).}
    \label{fig:bootstrap-ci}
  \end{figure}
\fi
Bootstrap intervals (Table~\ref{tab:bootstrap-ci}; visualized in Figure~\ref{fig:bootstrap-ci}) are tight for the Amazon low-rated proxy under the lexical baseline, suggesting that the aggregate category distribution is stable at this scale for our chosen taxonomy. Table~\ref{tab:bootstrap-ci-multi} extends this to $8$ method--dataset pairs in total: lexical on \DatasetAmazon, \DatasetTinder, and \DatasetGojek; TF-IDF on \DatasetAmazon{} and \DatasetAmazonAnn; and embedding (BGE-M3) on \DatasetAmazon, \DatasetTinder, and \DatasetGojek. Two observations follow: (i)~lexical and TF-IDF baselines on app-review proxies yield $W$ means rounding to $0$ with CIs $[0,0]$, which is "stable at zero" rather than "aligned"---i.e., these baselines deterministically fail to align under bootstrap resampling, consistent with the point estimates in Table~\ref{tab:alignment-main}; (ii)~$J_k$ on Tinder and Gojek shows mild but non-zero variance (CI width $\approx 0.18$) because top-$k$ set membership flickers near the boundary, while $W$ remains numerically zero, reinforcing that $J_k$ overstates alignment when the underlying mass overlap is negligible (the same effect documented in H2). The three embedding configurations are discussed under H3 below.
\IfFileExists{paper_ux/tables/bootstrap_ci_multi.tex}{\begin{table}[t]
\centering
\small
\caption{Bootstrap confidence intervals for alignment metrics across method--dataset pairs (200 iterations for lexical/TF-IDF, 50 for embedding; bootstrap resample size noted as $n$).}
\label{tab:bootstrap-ci-multi}
\resizebox{0.98\linewidth}{!}{%
\begin{tabular}{llrrrrr}
\toprule
Dataset & Method & $n$ & $J_k$ mean & $J_k$ 95\% CI & $W$ mean & $W$ 95\% CI \\
\midrule
\DatasetAmazon & \MethodLex & 79531 & 0.601 & [0.6, 0.6] & 0.0107 & [0.0104, 0.0112] \\
\DatasetAmazon & \MethodTFIDF & 10000 & 0.6 & [0.6, 0.6] & 0 & [0, 0] \\
\DatasetAmazonAnn & \MethodTFIDF & 8971 & 0.6 & [0.6, 0.6] & 0 & [0, 0] \\
\DatasetAmazon & \MethodEmbed & 200 & 0.804 & [0.778, 1] & 0 & [0, 0] \\
\DatasetTinder & \MethodLex & 10000 & 0.764 & [0.6, 0.778] & 0 & [0, 0] \\
\DatasetTinder & \MethodEmbed & 200 & 0.782 & [0.778, 0.778] & 0 & [0, 0] \\
\DatasetGojek & \MethodLex & 10000 & 0.74 & [0.6, 0.778] & 0 & [0, 0] \\
\DatasetGojek & \MethodEmbed & 200 & 0.774 & [0.778, 0.778] & 0 & [0, 0] \\
\bottomrule
\end{tabular}
}
\end{table}
}{}

\subsection{Failure-mode analysis: examples flagged as fabricated}
To make the grounding/fabrication proxies (Section~\ref{sec:limitations}) concrete, Table~\ref{tab:fabrication-examples} lists representative outputs that the adversarial judge flagged as fabricated, including the unsupported span and rationale. We note that the adversarial judge is intentionally one-sided---it emits items only when unsupported spans are detected---so every row is a positive case by construction; non-fabricated counterfactuals are not produced by this pass. This view is complementary to the aggregate hallucination/fabrication rates in Tables~\ref{tab:agent-ablation-hq} and~\ref{tab:agent-ablation-cq}: the aggregate rates are conservative lower bounds (stop-list induced), while the worked examples illustrate which kinds of friction phrasings the proxy actually catches.
\IfFileExists{paper_ux/tables/fabrication_examples.tex}{\begin{table}[t]
\centering
\footnotesize
\setlength{\tabcolsep}{4pt}
\caption{Top examples flagged as fabricated by an adversarial judge (4 representative cases, HQ ablation; the eval only emits items where unsupported spans are detected, so all rows are positive cases by construction).}
\label{tab:fabrication-examples}
\begin{tabular}{p{0.22\linewidth}rp{0.20\linewidth}p{0.42\linewidth}}
\toprule
Setting & Conf. & Friction & Unsupported span / rationale \\
\midrule
Hybrid (best-of-$N$ fallback, LLM judge) & 0.8 & navigation: threshold sliders unclear & span: Thresholds could be more intuitive. | Claims about UI elements (threshold sliders) are not supported by the UI snapshot text. \\
Best-of-$N$ (LLM judge) & 0.8 & navigation: unclear threshold sliders & span: friction-point: navigation: unclear threshold sliders | Claimed friction point and step 1 friction are supported by the UI snapshot, but the suggested fix is not. The UI snapshot shows thresh... \\
Best-of-$N$ (LLM judge) & 0.8 & navigation: falta de informacion sobre el sistema & span: falta de informacion sobre el sistema | The UI snapshot does not contain information about the suggested fixes or friction points. The sentiment analysis is also neutral, which contradicts th... \\
Hybrid (score-then-select, auto judge) & 0.8 & navigation: falta de contexto para entender a politica & span: falta de contexto para entender a politica | Claims about suggested fixes and friction points are not supported by the UI snapshot text. \\
\bottomrule
\end{tabular}
\end{table}
}{}

\subsection{Azure provider agent-ablation results}
Tables~\ref{tab:agent-ablation-hq} and \ref{tab:agent-ablation-cq} report the reduced ablation suites across three OpenAI deployments on Azure (\texttt{gpt-4.1}, \texttt{gpt-4.1-mini}, \texttt{gpt-5.2}).
\begin{itemize}
  \item \emph{Model sensitivity.} Under the same prompt/schema contract, \texttt{gpt-4.1} is strongest on quality (HQ hybrid-LLM quality $0.707$; CQ hybrid-auto and best-of-$N$ auto both at $0.726$), \texttt{gpt-4.1-mini} is consistently lower but still functional, and \texttt{gpt-5.2} collapses to near-degenerate behavior across every condition we ran (prefix rate $0$, grounding $0$, fabrication proxy $1.0$). The \texttt{gpt-5.2} failure is the only reproducible model-level breakdown we observed.
  \item \emph{Transient failure documented and re-measured.} An initial CQ batch for \texttt{gpt-4.1} (candidate temperature $0.7$, LLM judge) produced degenerate outputs (quality $0.4$, prefix $0$) on the hybrid and best-of-$N$ conditions. We re-executed those conditions with the identical configuration and \emph{the same model converged to non-degenerate behavior} (quality $0.71$--$0.73$, prefix $0.57$--$0.60$); the ablation tables report this latest measurement per condition. We therefore attribute the original CQ-v2 LLM-judge failures on \texttt{gpt-4.1} to a transient API/scheduling artifact, not a structural condition-specific failure of the prompt contract. The \texttt{gpt-5.2} collapse, in contrast, reproduced across batches and remains a real model-level failure.
\end{itemize}
The cost--quality frontier (Figure~\ref{fig:agent-ablation-pareto}) summarizes the resulting tradeoffs; auto-judge variants tend to dominate LLM-judge variants on cost without sacrificing quality, and the differences across models dwarf the differences across selection protocols. We report all conditions rather than filtering the \texttt{gpt-5.2} runs, because aggregate tables without model and condition labels would hide that breakdown.

\subsection{LLM label evaluation}
\IfFileExists{paper_ux/tables/appstore_eval.tex}{\begin{table}[t]
\centering
\small
\caption{LLM label evaluation on the Amazon Appstore annotated subset, using the same provider model family as the main UX agent-ablation matrix to keep the local-vs-provider policy consistent across tracks.}
\label{tab:appstore-eval}
\begin{tabular}{lrr}
\toprule
Model & Accuracy & Macro-F1 \\
\midrule
gpt-4.1 (n=500) & 0.556 & 0.438 \\
\bottomrule
\end{tabular}
\end{table}
}{}
On the annotated Amazon subset, \texttt{gpt-4.1} (Azure provider, same family as the main UX ablation) achieves moderate accuracy (Table~\ref{tab:appstore-eval}; accuracy $0.556$, macro-F1 $0.438$), consistent with uneven category difficulty and confusable labels (e.g., UX vs functionality). A cheap-baselines comparison (Table~\ref{tab:appstore-baselines}) sharpens the picture: an embedding + logistic-regression baseline (\texttt{nomic-embed-text}) attains higher accuracy ($0.635$) than the LLM, while the LLM still beats majority, lexical, and TF-IDF baselines on macro-F1. Table~\ref{tab:appstore-confusions} reports the most frequent label confusions and shows that the largest single error mode is \texttt{functionality\_features} being mapped to \texttt{ui\_ux}, which dominates the residual error budget.
Together, these results motivate our focus on distribution-level proxy alignment rather than treating LLM category assignments as reliable ground truth at the single-example level.
\IfFileExists{paper_ux/tables/appstore_baselines.tex}{\begin{table}[t]
\centering
\small
\caption{Cheap baselines for the annotated Appstore subset (same $n$ as Table~\ref{tab:appstore-eval}); embedding baseline uses a subset for cost.}
\label{tab:appstore-baselines}
\begin{tabular}{lrr}
\toprule
Baseline & Accuracy & Macro-F1 \\
\midrule
Majority & 0.432 & 0.086 \\
Lexical taxonomy match & 0.424 & 0.366 \\
TF-IDF + logistic regression (CV) & 0.492 & 0.146 \\
Embedding (nomic-embed-text) + logistic regression (CV) (n=200) & 0.635 & 0.336 \\
\bottomrule
\end{tabular}
\end{table}
}{}
\IfFileExists{paper_ux/tables/appstore_confusions.tex}{\begin{table}[t]
\centering
\small
\caption{Most frequent label confusions in LLM classification (Amazon annotated; n=500).}
\label{tab:appstore-confusions}
\resizebox{0.98\linewidth}{!}{%
\begin{tabular}{llr}
\toprule
Gold & Predicted & Count \\
\midrule
\texttt{performance\_stability} & \texttt{functionality\_features} & 30 \\
\texttt{functionality\_features} & \texttt{other} & 27 \\
\texttt{ui\_ux} & \texttt{functionality\_features} & 23 \\
\texttt{compatibility\_device} & \texttt{functionality\_features} & 20 \\
\texttt{support\_responsiveness} & \texttt{functionality\_features} & 17 \\
\texttt{functionality\_features} & \texttt{performance\_stability} & 16 \\
\texttt{functionality\_features} & \texttt{ui\_ux} & 12 \\
\texttt{compatibility\_device} & \texttt{performance\_stability} & 8 \\
\bottomrule
\end{tabular}
}
\end{table}
}{}

\subsection{Hypothesis evaluation}
We close with an explicit mapping from our hypotheses (Section~1) to results.
\textbf{H1} (embedding outperforms lexical under $W$) is \textbf{supported}: embedding yields higher $W$ on all primary app-review and support-tweet proxies
(e.g., Gojek: $W_\text{embed}=0.128$ vs $W_\text{lex}=0.000$; Amazon Low-rated: $W_\text{embed}=0.119$ vs $W_\text{lex}=0.009$).
The exception is Sentry OSS issues ($W_\text{lex}=0.138 > W_\text{embed}=0.072$), consistent with H1's ``tendency rather than universal'' framing.

\textbf{H2} ($J_k$ saturates and overstates alignment at large $k$) is \textbf{supported}: on the Amazon annotated proxy, lexical $J_k$ jumps from $0.333$ at $k=6$ to $1.0$ at $k=10$ while the weighted-Jaccard $W$ stays flat at $0.017$ across the same sensitivity sweep (canonical $W=0.013$ at $k=8$ in Table~\ref{tab:alignment-main}). $J_k$ is therefore misleading once $k$ approaches the taxonomy size.

\textbf{H3} (aggregate estimates are stable enough for iterative decision-making when reported with CIs) is \textbf{partially supported, with a systematic caveat about embedding-based $W$}. The multi-configuration bootstrap (Table~\ref{tab:bootstrap-ci-multi}; 200 iterations for lexical/TF-IDF, 50 for embedding) covers $8$ method--dataset pairs. Lexical and TF-IDF baselines on app-review proxies produce stable aggregate estimates: their $J_k$ CIs are zero-width or near-zero (top-$k$ set membership rarely changes under resampling), and their $W$ CIs are tight, often pinned at zero (deterministic non-alignment). For the lexical baseline on Amazon low-rated specifically, $W=0.0107$ with CI $[0.0104, 0.0112]$, indicating high stability at scale.

The embedding baseline (BGE-M3, 50 bootstrap iterations at $n=200$) tells a consistent and contrary story across all three proxies we ran. Its $J_k$ mean lands in $[0.77, 0.80]$ with tight CIs at $0.778$, but $W$ collapses to $0$ across all 50 iterations on \DatasetAmazon, \DatasetTinder, and \DatasetGojek alike---in sharp contrast to the single-shot point estimates of $W=0.119$, $0.087$, and $0.128$ respectively reported in Table~\ref{tab:alignment-main}. The three-way replication makes it implausible that the gap is a single-dataset anomaly: the embedding $W$ point estimates at this subsample size are \emph{systematically} unstable under bootstrap resampling. Practitioners should therefore report bootstrap CIs alongside any embedding-based alignment number before treating it as a target. In sum, H3 holds for lexical and TF-IDF baselines and for $J_k$ under embedding, but is \emph{contradicted} for embedding-based $W$ at this subsample size; characterizing the dependence on subsample size, anchor construction, and similarity threshold is left as future work.
